\documentclass[10pt, letterpaper]{article}

% Inhaltsverzeichnis für Pakettypen (nur für Übersicht im Header, wird nicht im Dokument angezeigt)
% 1. Seitenlayout und Ränder
% 2. Sprache und Zeichensatz
% 3. Mathematik und Theorem-Umgebungen
% 4. Eigene Makros
% 5. Diagramme und Grafiken
% 6. Tabellen und Aufzählungen
% 7. Inhaltsverzeichnis
% 8. Abschnittsüberschriften
% 9. Abstrakt-Umgebung
% 10. Todos/Notizen
% 11. Rahmen/Box-Umgebungen
% 12. Python-Integration
% 13. Literaturverwaltung
% 14. Hyperlinks
% 15. Absatzeinstellungen
% 16. Umgebungen
% 17  Graphik
% 18  Extra
% 00. Titel und Autor

\usepackage{slashed}

% --- 1. Seitenlayout und Ränder ---
\usepackage[margin=3cm]{geometry}

% --- 2. Sprache und Zeichensatz ---
\usepackage[english]{babel}
\usepackage[T1]{fontenc}
\usepackage[utf8]{inputenc}

% --- 3. Mathematik und Theorem-Umgebungen ---
\usepackage{amsmath, amssymb, amsthm}
\usepackage{mathrsfs}
\DeclareMathOperator{\WF}{WF}

% --- 4. Eigene Makros ---
\usepackage{xcolor}
\newcommand{\SKP}{\langle\cdot,\cdot\rangle}
\newcommand{\id}{\operatorname{id}}
\newcommand{\sgn}{\operatorname{sgn}}
\newcommand{\Ad}{\operatorname{Ad}}
\newcommand{\ad}{\operatorname{ad}}
\newcommand{\K}{\mathbb{K}}
\newcommand{\R}{\mathbb{R}}
\newcommand{\N}{\mathbb{N}}
\newcommand{\Q}{\mathbb{Q}}
\newcommand{\Z}{\mathbb{Z}}
\newcommand{\C}{\mathbb{C}}
\newcommand{\QU}{\mathbb{H}}
\newcommand{\Cinf}{C^{\infty}}
\newcommand{\Mod}{\mathrm{Mod}}
\newcommand{\Alg}{\mathrm{Alg}}
\newcommand{\Diff}{\mathrm{Diff}}
\newcommand{\Iso}{\mathrm{Iso}}
\newcommand{\Aut}{\mathrm{Aut}}
\newcommand{\End}{\mathrm{End}}
\newcommand{\Hom}{\mathrm{Hom}}
\newcommand{\Cl}{\mathrm{Cl}}
\newcommand{\entwurf}[1]{\textcolor{red}{#1}}
\newcommand{\GL}{\mathrm{GL}}
\newcommand{\Mat}{\mathrm{Mat}}
\newcommand{\SL}{\mathrm{SL}}
\newcommand{\SO}{\mathrm{SO}}
\newcommand{\SU}{\mathrm{SU}}
\newcommand{\Sp}{\mathrm{Sp}}
\newcommand{\Spin}{\mathrm{Spin}}
\newcommand{\Pin}{\mathrm{Pin}}
\newcommand{\Lor}{\mathrm{Lor}}
\newcommand{\Ogroup}{\mathrm{O}}
\newcommand{\U}{\mathrm{U}}
\newcommand{\CP}{\mathbb{C} P}
\newcommand{\RP}{\mathbb{R} P}

% --- 5. Diagramme und Grafiken ---
\usepackage{graphicx}
\graphicspath{{../../Library/Images/}}
\usepackage[export]{adjustbox}
\usepackage{tikz}
\usetikzlibrary{decorations.pathreplacing, arrows.meta, positioning}
\usepackage{tikz-cd}

% --- 6. Tabellen und Aufzählungen ---
\usepackage{enumitem}
\setlist[itemize]{left=0.5cm}

\newenvironment{romanenum}[1][]
  {%
    \ifx&#1&
    \else
      \textbf{#1}\quad
    \fi
    \begin{enumerate}[label=\roman*)]
  }
  {%
    \end{enumerate}%
  }

% --- 7. Inhaltsverzeichnis ---
\usepackage{tocloft}
\renewcommand{\cftsecfont}{\footnotesize}
\renewcommand{\cftsubsecfont}{\footnotesize}
\renewcommand{\cftsubsubsecfont}{\footnotesize}
\renewcommand{\cftsecpagefont}{\footnotesize}
\renewcommand{\cftsubsecpagefont}{\footnotesize}
\renewcommand{\cftsubsubsecpagefont}{\footnotesize}
\usepackage{etoc}

% --- 8. Abschnittsüberschriften ---
\usepackage{titlesec}
\titleformat{\section}{\normalfont\large\bfseries}{\thesection}{1em}{}
\titleformat{\subsection}{\normalfont\normalsize\bfseries}{\thesubsection}{0.5em}{}
\titleformat{\subsubsection}{\normalfont\normalsize\bfseries}{\thesubsubsection}{0.5em}{}
\setcounter{secnumdepth}{4}

% --- 9. Abstrakt-Umgebung ---
\usepackage{changepage}
\renewenvironment{abstract}
  {
    \begin{adjustwidth}{1.5cm}{1.5cm}
    \small
    \textsc{Abstract. –}%
  }
  {
    \end{adjustwidth}
  }

% --- 10. Todos/Notizen ---
\usepackage{todonotes}
% Hellblau definieren
\definecolor{hellblau}{rgb}{0.3,0.6,1}
\newenvironment{note}{\par\begingroup\color{hellblau}\itshape}{\par\endgroup}

% --- 11. Rahmen/Box-Umgebungen ---
\usepackage{mdframed}
\usepackage{tcolorbox}
\colorlet{shadecolor}{gray!25}

\newenvironment{customTheorem}
  {\vspace{10pt}%
   \begin{mdframed}[
     backgroundcolor=gray!20,
     linewidth=0pt,
     innertopmargin=10pt,
     innerbottommargin=10pt,
     skipabove=\dimexpr\topsep+\ht\strutbox\relax,
     skipbelow=\topsep,
   ]}
  {\end{mdframed}
   \vspace{10pt}%
  }

% --- 12. Python-Integration ---
% (Deaktiviert in dieser Version, aktiviere bei Bedarf)
% \usepackage{pythontex}
% \usepackage[makestderr]{pythontex}

% --- 13. Literaturverwaltung ---
\usepackage{csquotes}
\usepackage[backend=biber, style=alphabetic, citestyle=alphabetic]{biblatex}
\addbibresource{bibliography.bib}

% --- 14. Hyperlinks ---
\usepackage{hyperref}
\hypersetup{
  colorlinks   = true,
  urlcolor     = blue,
  linkcolor    = blue,
  citecolor    = blue,
  frenchlinks  = true
}

% --- 15. Absatzeinstellungen ---
\usepackage[parfill]{parskip}
\sloppy

% --- 16. Umgebungen ---
\usepackage{thmtools}

\newcommand{\CustomHeading}[3]{%
  \par\medskip\noindent%
  \textbf{#1 #2} \textnormal{(#3)}.\enskip%
}

\newenvironment{DEF}[2]{\CustomHeading{Definition}{#1}{#2}}{}
\newenvironment{PROP}[2]{\CustomHeading{Proposition}{#1}{#2}}{}
\newenvironment{THEO}[2]{\CustomHeading{Theorem}{#1}{#2}}{}
\newenvironment{LEM}[2]{\CustomHeading{Lemma}{#1}{#2}}{}
\newenvironment{KORO}[2]{\CustomHeading{Corollar}{#1}{#2}}{}
\newenvironment{REM}[2]{\CustomHeading{Remark}{#1}{#2}}{}
\newenvironment{EXA}[2]{\CustomHeading{Example}{#1}{#2}}{}
\newenvironment{STUD}[2]{\CustomHeading{Study}{#1}{#2}}{}
\newenvironment{CONC}[2]{\CustomHeading{Concept}{#1}{#2}}{}
\newenvironment{OTH}[2]{\CustomHeading{Other}{#1}{#2}}{}
\newenvironment{EXE}[2]{\CustomHeading{Exercise}{#1}{#2}}{}
\newenvironment{MOT}[2]{\CustomHeading{Motivation}{#1}{#2}}{}
\newenvironment{PROOF}[2]{\CustomHeading{Proof}{#1}{#2}}{}

% --- Unit Umgebung für Source-Inhalte ---
\usepackage{mdframed}
\newmdenv[
  linewidth=1pt,
  topline=false,
  bottomline=false,
  rightline=false,
  leftmargin=0cm,
  rightmargin=0cm,
  skipabove=10pt,
  skipbelow=10pt,
  innertopmargin=0.5\baselineskip,
  innerbottommargin=0.5\baselineskip,
  backgroundcolor=gray!10,
  linecolor=gray
]{unitbox}

\newenvironment{unit}[1]
  {\begin{unitbox}\textbf{Unit #1}\par\smallskip}
  {\end{unitbox}}

% --- 17. ... ---

% --- 18. Extras ---
\usepackage{stmaryrd}
\usepackage{bbold}  % falls du \mathbb{1} nutzen willst

% --- 19. Inhaltsverzeichnis ---
\usepackage{tocloft}

% Abstand zwischen Einträgen
\setlength{\cftbeforesecskip}{2pt}      % Sections
\setlength{\cftbeforesubsecskip}{1pt}   % Subsections
\setlength{\cftbeforesubsubsecskip}{0pt}
\renewcommand{\cfttoctitlefont}{\large\bfseries}

% --- 00. Titel und Autor ---
\title{Matlab Dictionary}
\author{Tim Jaschik}
\date{\today}

\begin{document}

\maketitle
\rule{\textwidth}{0.5pt}
\begin{abstract}
Kurze Beschreibung …
\end{abstract}
\rule{\textwidth}{0.5pt}
\vspace{0.5cm}

\tableofcontents

\pagebreak

\section{Systematische Übersicht zu \texttt{fprintf} in MATLAB}

\subsection{Grundidee}

Der Befehl \texttt{fprintf} dient in MATLAB dazu, formatierte Ausgaben im Command Window zu erzeugen. Die Grundstruktur lautet:

\begin{verbatim}
fprintf(formatString, value1, value2, ...)
\end{verbatim}

Dabei ist \texttt{formatString} eine Zeichenkette, die sowohl normalen Text als auch sogenannte Platzhalter enthält. Die Werte \texttt{value1}, \texttt{value2}, \dots werden der Reihe nach in diese Platzhalter eingesetzt.

Ein einfaches Beispiel ist:

\begin{verbatim}
x = 3.14159;
fprintf('Der Wert von x ist %.2f\n', x);
\end{verbatim}

Die Ausgabe ist:

\begin{verbatim}
Der Wert von x ist 3.14
\end{verbatim}

Der Ausdruck \texttt{\%.2f} bedeutet: Setze hier eine Dezimalzahl mit zwei Nachkommastellen ein. Der Ausdruck \texttt{\textbackslash n} erzeugt einen Zeilenumbruch.

\subsection{Allgemeine Struktur eines Platzhalters}

Ein Platzhalter hat typischerweise die Form

\[
\texttt{\%[flags][width][.precision]type}.
\]

Die wichtigsten Bestandteile sind:

\begin{enumerate}
    \item \texttt{\%}: Beginn eines Platzhalters.
    \item \texttt{width}: minimale Gesamtbreite der Ausgabe.
    \item \texttt{.precision}: Anzahl der Nachkommastellen oder signifikanten Stellen.
    \item \texttt{type}: Datentyp beziehungsweise Ausgabeformat.
\end{enumerate}

Beispiel:

\begin{verbatim}
fprintf('%8.2f\n', 3.14159);
\end{verbatim}

Hier bedeutet:

\[
\texttt{\%8.2f}
\]

\begin{enumerate}
    \item \texttt{8}: mindestens acht Zeichen Gesamtbreite,
    \item \texttt{.2}: zwei Nachkommastellen,
    \item \texttt{f}: Ausgabe als Dezimalzahl.
\end{enumerate}

Die Ausgabe ist ungefähr:

\begin{verbatim}
    3.14
\end{verbatim}

\subsection{Wichtige Formattypen}

\begin{center}
\begin{tabular}{lll}
\hline
Platzhalter & Bedeutung & Beispielausgabe \\
\hline
\texttt{\%d} & Ganze Zahl & \texttt{42} \\
\texttt{\%i} & Ganze Zahl & \texttt{42} \\
\texttt{\%f} & Dezimalzahl & \texttt{3.141593} \\
\texttt{\%.2f} & Dezimalzahl mit 2 Nachkommastellen & \texttt{3.14} \\
\texttt{\%e} & Wissenschaftliche Schreibweise & \texttt{3.141593e+00} \\
\texttt{\%.2e} & Wissenschaftliche Schreibweise mit 2 Nachkommastellen & \texttt{3.14e+00} \\
\texttt{\%g} & Kompakte Darstellung & \texttt{3.14159} \\
\texttt{\%s} & Zeichenkette/String & \texttt{Hallo} \\
\texttt{\%c} & Einzelnes Zeichen & \texttt{A} \\
\texttt{\%\%} & Prozentzeichen & \texttt{\%} \\
\hline
\end{tabular}
\end{center}

\subsection{Ausgabe ganzer Zahlen}

Für ganze Zahlen verwendet man meistens \texttt{\%d}.

\begin{verbatim}
k = 17;
fprintf('Die Anzahl der Iterationen ist %d.\n', k);
\end{verbatim}

Ausgabe:

\begin{verbatim}
Die Anzahl der Iterationen ist 17.
\end{verbatim}

Typische Anwendung bei numerischen Algorithmen:

\begin{verbatim}
fprintf('k_GS = %d\n', k_GS);
\end{verbatim}

\subsection{Ausgabe von Dezimalzahlen}

Für Dezimalzahlen verwendet man \texttt{\%f}. Standardmäßig werden sechs Nachkommastellen angezeigt.

\begin{verbatim}
x = 3.1415926535;
fprintf('x = %f\n', x);
\end{verbatim}

Ausgabe:

\begin{verbatim}
x = 3.141593
\end{verbatim}

Möchte man die Anzahl der Nachkommastellen kontrollieren, verwendet man zum Beispiel:

\begin{verbatim}
fprintf('x = %.2f\n', x);
fprintf('x = %.6f\n', x);
fprintf('x = %.10f\n', x);
\end{verbatim}

Ausgabe:

\begin{verbatim}
x = 3.14
x = 3.141593
x = 3.1415926535
\end{verbatim}

\subsection{Wissenschaftliche Schreibweise}

Für sehr kleine oder sehr große Zahlen ist \texttt{\%e} oft besser geeignet.

\begin{verbatim}
err = 0.00000012345;
fprintf('error = %.4e\n', err);
\end{verbatim}

Ausgabe:

\begin{verbatim}
error = 1.2345e-07
\end{verbatim}

Das ist besonders nützlich für Fehler, Toleranzen und Residuen:

\begin{verbatim}
fprintf('residual norm = %.4e\n', norm(A*x-b));
\end{verbatim}

\subsection{Kompakte numerische Darstellung mit \texttt{\%g}}

Der Platzhalter \texttt{\%g} wählt automatisch eine kompakte Darstellung. MATLAB entscheidet also je nach Zahl, ob eine normale Dezimaldarstellung oder wissenschaftliche Schreibweise sinnvoller ist.

\begin{verbatim}
fprintf('x = %g\n', 3.1415926535);
fprintf('x = %g\n', 0.00000012345);
\end{verbatim}

Mögliche Ausgabe:

\begin{verbatim}
x = 3.14159
x = 1.2345e-07
\end{verbatim}

\subsection{Ausgabe von Strings}

Für Textvariablen verwendet man \texttt{\%s}.

\begin{verbatim}
method = 'Gauss-Seidel';
fprintf('Method: %s\n', method);
\end{verbatim}

Ausgabe:

\begin{verbatim}
Method: Gauss-Seidel
\end{verbatim}

Mehrere Werte werden der Reihe nach eingesetzt:

\begin{verbatim}
method = 'Gauss-Seidel';
k = 23;
fprintf('%s converged after %d iterations.\n', method, k);
\end{verbatim}

Ausgabe:

\begin{verbatim}
Gauss-Seidel converged after 23 iterations.
\end{verbatim}

\subsection{Zeilenumbrüche und Sonderzeichen}

Wichtige Steuerzeichen sind:

\begin{center}
\begin{tabular}{ll}
\hline
Code & Bedeutung \\
\hline
\texttt{\textbackslash n} & Zeilenumbruch \\
\texttt{\textbackslash t} & Tabulator \\
\texttt{\textbackslash\textbackslash} & Backslash ausgeben \\
\texttt{\%\%} & Prozentzeichen ausgeben \\
\hline
\end{tabular}
\end{center}

Beispiel:

\begin{verbatim}
fprintf('Name:\tGauss-Seidel\n');
fprintf('Tolerance:\t%.2e\n', 1e-8);
\end{verbatim}

Ausgabe:

\begin{verbatim}
Name:       Gauss-Seidel
Tolerance:  1.00e-08
\end{verbatim}

\subsection{Breite und Ausrichtung}

Man kann festlegen, wie viel Platz eine Zahl mindestens einnehmen soll. Das ist besonders nützlich für tabellarische Ausgaben.

\begin{verbatim}
fprintf('%10.4f\n', 3.14159);
fprintf('%10.4f\n', 123.456);
\end{verbatim}

Ausgabe:

\begin{verbatim}
    3.1416
  123.4560
\end{verbatim}

Hier bedeutet:

\[
\texttt{\%10.4f}
\]

\begin{enumerate}
    \item \texttt{10}: mindestens zehn Zeichen Breite,
    \item \texttt{.4}: vier Nachkommastellen,
    \item \texttt{f}: Dezimalzahl.
\end{enumerate}

Linksbündige Ausgabe erhält man mit einem Minuszeichen:

\begin{verbatim}
fprintf('%-10.4f Ende\n', 3.14159);
\end{verbatim}

Ausgabe:

\begin{verbatim}
3.1416     Ende
\end{verbatim}

\subsection{Ausgabe von Vektoren}

Wenn man einen Vektor direkt mit \texttt{fprintf} ausgibt, wird derselbe Platzhalter auf alle Einträge angewendet.

\begin{verbatim}
x = [1.23456; 2.34567; 3.45678];

fprintf('x = [');
fprintf(' %.4f', x);
fprintf(' ]\n');
\end{verbatim}

Ausgabe:

\begin{verbatim}
x = [ 1.2346 2.3457 3.4568 ]
\end{verbatim}

Für einen Spaltenvektor mit Komponentenbeschriftung:

\begin{verbatim}
for j = 1:length(x)
    fprintf('x_%d = %.6f\n', j, x(j));
end
\end{verbatim}

Ausgabe:

\begin{verbatim}
x_1 = 1.234560
x_2 = 2.345670
x_3 = 3.456780
\end{verbatim}

\subsection{Tabellarische Ausgabe}

Für Ergebnisse verschiedener Algorithmen kann man eine einfache Tabelle im Command Window erzeugen.

\begin{verbatim}
fprintf('\n%-15s %10s %15s\n', 'Method', 'Iter', 'Residual');
fprintf('%-15s %10d %15.4e\n', 'Jacobi', k_GJ, norm(A*x_GJ-b));
fprintf('%-15s %10d %15.4e\n', 'Gauss-Seidel', k_GS, norm(A*x_GS-b));
\end{verbatim}

Mögliche Ausgabe:

\begin{verbatim}
Method                Iter        Residual
Jacobi                  42      3.2154e-09
Gauss-Seidel            18      7.4512e-10
\end{verbatim}

Dabei bedeutet:

\begin{enumerate}
    \item \texttt{\%-15s}: String, linksbündig, Breite 15.
    \item \texttt{\%10d}: Ganze Zahl, rechtsbündig, Breite 10.
    \item \texttt{\%15.4e}: Wissenschaftliche Zahl, Breite 15, vier Nachkommastellen.
\end{enumerate}

\subsection{Typische Formatierungen für numerische Algorithmen}

Für Iterationszahlen:

\begin{verbatim}
fprintf('iterations = %d\n', k);
\end{verbatim}

Für Lösungen mit sechs Nachkommastellen:

\begin{verbatim}
fprintf('x = %.6f\n', x);
\end{verbatim}

Für Fehler oder Residuen:

\begin{verbatim}
fprintf('error = %.4e\n', err);
fprintf('residual = %.4e\n', norm(A*x-b));
\end{verbatim}

Für kompakte Parameterwerte:

\begin{verbatim}
fprintf('lambda = %.2f\n', lambda);
\end{verbatim}

Für eine kombinierte Ausgabe:

\begin{verbatim}
fprintf('lambda = %.2f, iterations = %d, residual = %.4e\n', ...
    lambda, k, norm(A*x-b));
\end{verbatim}

\subsection{Häufige Fehler}

\subsubsection{Falscher Platzhalter}

Für Dezimalzahlen sollte man nicht \texttt{\%d}, sondern \texttt{\%f}, \texttt{\%e} oder \texttt{\%g} verwenden.

Ungünstig:

\begin{verbatim}
fprintf('x = %d\n', 3.14159);
\end{verbatim}

Besser:

\begin{verbatim}
fprintf('x = %.6f\n', 3.14159);
\end{verbatim}

\subsubsection{Zu wenige oder zu viele Werte}

Die Anzahl der Platzhalter sollte zur Anzahl der eingefügten Werte passen.

\begin{verbatim}
fprintf('x = %.2f, y = %.2f\n', x, y);
\end{verbatim}

Hier gibt es zwei Platzhalter und zwei Werte.

\subsubsection{Prozentzeichen vergessen zu escapen}

Ungünstig:

\begin{verbatim}
fprintf('Convergence rate = 95%\n');
\end{verbatim}

Besser:

\begin{verbatim}
fprintf('Convergence rate = 95%%\n');
\end{verbatim}

\subsection{Merkschema}

Für die meisten numerischen Anwendungen reichen die folgenden Muster:

\begin{center}
\begin{tabular}{ll}
\hline
Zweck & MATLAB-Code \\
\hline
Ganze Zahl & \texttt{fprintf('k = \%d\textbackslash n', k)} \\
Dezimalzahl & \texttt{fprintf('x = \%.6f\textbackslash n', x)} \\
Kleine/große Zahl & \texttt{fprintf('err = \%.4e\textbackslash n', err)} \\
String & \texttt{fprintf('Method = \%s\textbackslash n', method)} \\
Prozentzeichen & \texttt{fprintf('rate = 95\%\%\textbackslash n')} \\
Vektor in einer Zeile & \texttt{fprintf(' \%.6f', x)} \\
\hline
\end{tabular}
\end{center}

\pagebreak
\printbibliography
\end{document}